By Theorem \claimref{synthesis.quarter-from-temperature}, under (P, \claimref{singularity.p-univalence}) and (R, \claimref{singularity.r-regularity-fixed-locus}) the hole's thermodynamic entropy is $S = A/4$. The hypothesis $\sum_{j} s(x_{j}) = A/4$ gives $|J|\,\langle s\rangle = A/4$, so $\kappa_{c} = |J|/A = 1/(4\langle s\rangle)$, which is (i).

By Theorem \claimref{synthesis.boundary-mode-thermal-entropy}, $s(x) \le \ln 2$ with equality exactly at $x = 0$, so $\langle s\rangle \le \ln 2$ with equality exactly when every $x_{j} = 0$, that is every $\omega_{j} = 0$. Hence $\kappa_{c} = 1/(4\langle s\rangle) \ge 1/(4\ln 2)$, with the same case of equality, which is (ii). In the capacity convention T08 sets $\kappa_{c} = 1/(4\ln 2)$ and the boundary-mode prediction states $\nu = 1/(16\pi\ln 2)$; by (ii) these hold under the hypothesis exactly when the counted modes are at zero weight.

For (iii): T08's remark on $\nu$ already identifies $\ln 2$ as the entropy of a maximally mixed fermionic mode. At equality $\kappa_{c} = 1/(4\ln 2)$ with $\ln 2 = s(0)$, the entropy of the state $\tfrac12\mathbf 1$ of one mode, so the transcendental factor is a property of the state of a mode and not of a tally; a rule that assigns a rational number of modes per cell of rational area produces a residue in $\mathbb{Q}\pi^{k}$, which T08 shows never equals $1/(4\ln 2)$ for $k = 0$.

For (iv): under (P) and (R$_{\Lambda}$, \claimref{synthesis.r-lambda-regularity-cosmological-horizon}) each mode of the cosmological horizon is Fermi-Dirac at $T_{\Lambda} = \kappa_{\Lambda}/2\pi$ (Theorem \claimref{synthesis.de-sitter-temperature-frozen-point-norm}); the computation of Theorem \claimref{synthesis.boundary-mode-thermal-entropy} uses only the Fermi-Dirac form, so with $x_{j} = 2\pi\omega_{j}/\kappa_{\Lambda}$ the same bound $s \le \ln 2$ holds, with equality exactly at zero weight. With $S_{\Lambda} = \pi N^{2}$ (Theorem \claimref{synthesis.de-sitter-entropy-pi-n-squared}) the hypothesis $\sum_{j} s(x_{j}) = S_{\Lambda}$ gives $|J| = \pi N^{2}/\langle s\rangle \ge \pi N^{2}/\ln 2$, with equality exactly at zero weight; since Definition \claimref{obj.synthesis.cosmological-horizon-boundary-modes} has $\omega_{j} > 0$, the inequality is strict.

For (v): with $n = 1/(e^{x} + 1)$, $s(x) = \ln(1 + e^{-x}) + xn > xn = \beta\omega n$ for every finite $x$, since $\ln(1 + e^{-x}) > 0$. Summing, $\beta E = \sum_{j}\beta\omega_{j}n_{j} < \sum_{j} s(x_{j}) = A/4 = S$, so $E < TS$, and $TS = (1/8\pi M)(4\pi M^{2}) = M/2$. At equality in (ii) every $\omega_{j} = 0$ and $E = 0$.

For (vi): under (Q) (\claimref{singularity.q-joint-quasi-freeness-s-wave} records the weights $\omega_{k} = \kappa k$, the sea level at $k = 0$) an array's modes have $x = 2\pi k$, $k \in \mathbb{Z}$, and since $s$ is even its entropy is $s(0) + 2\sum_{k \ge 1} s(2\pi k) = \ln 2 + \varepsilon$. The check brackets $\varepsilon$ with rational bounds: the terms $k = 1, 2, 3$ by rational enclosures of $\pi$ and the exponential, and the tail by $s(x) < (1 + x)e^{-x}$ for $x > 0$. The three readings follow: every mode counted gives $\langle s\rangle = (\ln 2 + \varepsilon)/m$ for an array of $m$ modes, which tends to $0$ as $m$ grows; arrays counted give $\kappa_{c} = 1/(4(\ln 2 + \varepsilon))$; only zero-weight modes counted give $\sum s = |J|\ln 2$ while the arrays hold $|J|(\ln 2 + \varepsilon)$.
